\documentclass[12pt]{article}
\usepackage[utf8]{inputenc}
\usepackage[T1]{fontenc}
\usepackage{hyperref}
\title{Vinitiri Language Resources for K-12 Curriculum Development}
\author{Compiled by Claude Sonnet 5 (claude-sonnet-5) for Andr\'as Kornai}
\date{}
\begin{document}
\maketitle
\section{Introduction}
Vinitiri (ISO 639-3 vmg, Glottocode mini1251) is an Austronesian language of Papua New Guinea. The search results describe it as spoken in East New Britain, mainly around the villages of Kabaira and Vunamarita. They give roughly 2,000 speakers, which differs from the figure of about 200 in the task description. The search results also include Vitu, a different language of West New Britain, which must not be confused with Vinitiri. Very little exists beyond one grammar.
\section{Primary Educational Resources}
None found.
\section{Linguistic and Reference Materials}
\begin{itemize}
\item Sheena Van Der Mark, \emph{A Grammar of Vinitiri: An Austronesian language from Papua New Guinea}, doctoral dissertation, La Trobe University, 2007 (as cited in Grambank): \url{https://grambank.clld.org/sources/gVandermarkVinitiri}
\item Grambank language entry: \url{https://grambank.clld.org/languages/mini1251}
\end{itemize}
\section{Religious and Community Resources}
None found.
\section{Digital and Audio Resources}
None found.
\section{Curriculum Development Context}
Documentation is limited to a single reference grammar. No pedagogical materials, orthography guides or school programmes were found. Any curriculum would need to start from that grammar together with work with the community.
\end{document}